\textbf{``Preventing Wrong Decisions in Smart Building Systems''}

\emph{\textbf{Minor Project Report}}

\emph{submitted in partial fulfillment of the}

\emph{requirements for the award of the degree of}

\textbf{BACHELOR OF TECHNOLOGY}

in

\textbf{COMPUTER SCIENCE \& TECHNOLOGY}

\textbf{by}

{\def\LTcaptype{none} % do not increment counter
\begin{longtable}[]{@{}
  >{\centering\arraybackslash}p{(\linewidth - 2\tabcolsep) * \real{0.2724}}
  >{\centering\arraybackslash}p{(\linewidth - 2\tabcolsep) * \real{0.3320}}@{}}
\toprule\noalign{}
\begin{minipage}[b]{\linewidth}\centering
\textbf{Name}
\end{minipage} & \begin{minipage}[b]{\linewidth}\centering
\textbf{Roll No.}
\end{minipage} \\
\midrule\noalign{}
\endhead
\bottomrule\noalign{}
\endlastfoot
Rishi Kumar & 500126773 \\
Pranav VR & 500119049 \\
Aman Rajpoot & 500123398 \\
Amalkrishnan Ajay & 500120559 \\
\end{longtable}
}

\emph{\textbf{Under the guidance of}}

Dr. Amit Gurung

\includegraphics[width=2.18889in,height=1.44167in,alt={Research}]{media/image1.jpeg}

\textbf{School of Computer Science, UPES}

Bidholi, Via Prem Nagar, Dehradun, Uttarakhand

May -- 2026

\textbf{CANDIDATE'S DECLARATION}

We hereby certify that the project work entitled \textbf{``Preventing
Wrong Decisions in Smart Building Systems''} in partial fulfilment of
the requirements for the award of the Degree of BACHELOR OF TECHNOLOGY
in COMPUTER SCIENCE AND ENGINEERING with specialization in Cybersecurity
and Digital Forensic and submitted to the Department of Systemics,
School of Computer Science, University of Petroleum \& Energy Studies,
Dehradun, is an authentic record of my our work carried out during a
period from \textbf{Jan,} 2026 to \textbf{May}, \textbf{2026} under the
supervision of \textbf{Amit Gurung, \_\_\_\_Cluster, UPES}.

The matter presented in this project has not been submitted by us for
the award of any other degree of this or any other University.

\textbf{Rishi Kumar Pranav VR Aman Rajpoot Amalkrishnan Ajay}

500126773 500119049 500123398 500120559

This is to certify that the above statement made by the candidate is
correct to the best of my knowledge.

Date: \_\_\_\_\_\_\_\_\_ 2026 Project Guide

\textbf{Acknowledgement}

We sincerely thanks to our respected Head Department, \textbf{Mr. Hitesh
Kumar} for his great support in doing our project.

We sincerely express our gratitude to \textbf{Dr. Srinivasan
Ramachandran}, Faculty, School of Computer Science (AI and Computer
Vision Cluster), UPES, for his constant guidance and support. We
especially thank him for providing the necessary hardware resources and
for his valuable mentorship in helping us understand and implement IoT
technologies throughout this project.

We are also grateful to Dean SoCS UPES for giving us the necessary
facilities to carry out our project work successfully. We also thanks to
our Course Coordinator, \textbf{Azra Nazir} and our Activity Coordinator
\textbf{Nitesh Kumar Singh} for providing timely support and information
during the completion of this project.

We would like to thank all our friends for their help and constructive
criticism during our project work. Finally, we have no words to express
our sincere gratitude to our parents who have shown us this world and
for every support they have given us.

{\def\LTcaptype{none} % do not increment counter
\begin{longtable}[]{@{}
  >{\raggedright\arraybackslash}p{(\linewidth - 8\tabcolsep) * \real{0.1590}}
  >{\raggedright\arraybackslash}p{(\linewidth - 8\tabcolsep) * \real{0.2029}}
  >{\raggedright\arraybackslash}p{(\linewidth - 8\tabcolsep) * \real{0.2029}}
  >{\raggedright\arraybackslash}p{(\linewidth - 8\tabcolsep) * \real{0.2029}}
  >{\raggedright\arraybackslash}p{(\linewidth - 8\tabcolsep) * \real{0.2323}}@{}}
\toprule\noalign{}
\begin{minipage}[b]{\linewidth}\raggedright
\textbf{Name}
\end{minipage} & \begin{minipage}[b]{\linewidth}\raggedright
\textbf{Rishi Kumar}
\end{minipage} & \begin{minipage}[b]{\linewidth}\raggedright
\textbf{Pranav VR}
\end{minipage} & \begin{minipage}[b]{\linewidth}\raggedright
\textbf{Aman Rajpoot}
\end{minipage} & \begin{minipage}[b]{\linewidth}\raggedright
\textbf{Amalkrishnan Ajay}
\end{minipage} \\
\midrule\noalign{}
\endhead
\bottomrule\noalign{}
\endlastfoot
SAP Id. & 500126773 & 500119049 & 500123398 & 500120559 \\
\end{longtable}
}

\textbf{Abstract}

The rapid growth of smart building systems and IoT-based automation has
increased reliance on sensor-driven decision-making processes. Most
existing systems focus on securing data transmission using protocols
such as MQTT with TLS, ensuring that data is protected during
communication. However, these systems assume that incoming sensor data
is always correct. In reality, sensors can be faulty or manipulated
while still sending believable values, leading to incorrect decisions
without any warning. This can result in energy inefficiency, system
malfunction, and potential safety risks.

To address this issue, the proposed Hybrid IoT Intrusion Detection
System (IDS) introduces a mechanism to verify both the security and
correctness of data. The system collects real-time data from multiple
IoT sensor nodes connected through ESP32 devices and an MQTT broker. A
real-world attack scenario was also performed using ARP spoofing,
allowing live interception and manipulation of sensor data. This helped
in generating realistic datasets for training and testing the system.

The system combines multiple detection approaches, including rule-based
logic and machine learning, to identify abnormal patterns in sensor data
over time. Instead of relying on single readings, it analyses sequences
of data to detect subtle changes such as sudden spikes, gradual drift,
repeated patterns, or unexpected drops. In addition, a network
monitoring layer observes traffic behaviour to detect suspicious
activities before they affect the system.

All these components work together in a hybrid decision process,
improving the accuracy and reliability of anomaly detection. The system
also includes real-time visualization to help monitor sensor behaviour
and identify issues easily.

Overall, the proposed solution ensures that automated decisions are
based on trustworthy data rather than blindly accepted inputs. It
provides a practical and scalable approach to improving the reliability
and security of smart building systems and other IoT-based environments.

\textbf{TABLE OF CONTENTS}

\textbf{S.No. Contents Page No}

\begin{enumerate}
\def\labelenumi{\arabic{enumi}.}
\item
  \textbf{Introduction 1}

  \begin{enumerate}
  \def\labelenumii{\arabic{enumii}.}
  \item
    History 2
  \item
    Requirement Analysis
  \end{enumerate}
\end{enumerate}

\begin{quote}
1.2.1 Functional Requirements

1.2.2 Non-Functional Requirements
\end{quote}

\begin{enumerate}
\def\labelenumi{\arabic{enumi}.}
\setcounter{enumi}{2}
\item
  Main Objective 4
\item
  Sub Objectives 4
\item
  System Workflow 5
\end{enumerate}

\begin{enumerate}
\def\labelenumi{\arabic{enumi}.}
\setcounter{enumi}{1}
\item
  \textbf{System Analysis 6}

  \begin{enumerate}
  \def\labelenumii{\arabic{enumii}.}
  \item
    Introduction 6
  \item
    Existing System 7
  \item
    Comparative Analysis 7
  \item
    Research Gap 9
  \end{enumerate}
\item
  \textbf{System Design 10}
\end{enumerate}

3.1 Proposed Solution 10

3.2 System Architecture 11

3.3 Modules

\begin{quote}
3.4 Design Diagrams 54

3.4.1 Use Case Diagram 43

3.4.2 Activity Diagram 54
\end{quote}

\begin{enumerate}
\def\labelenumi{\arabic{enumi}.}
\setcounter{enumi}{3}
\item
  \textbf{Technology Stack and System Specification 18}

  \begin{enumerate}
  \def\labelenumii{\arabic{enumii}.}
  \item
    Overview 18
  \end{enumerate}
\end{enumerate}

4.2 Hardware Components 19

\begin{quote}
4.3 Software Components 20

4.4 System Configuration 21

4.5 Design Considerations 23

4.5.1 Scalability 32

4.5.2 Reliability 32

4.5.3 Security 23
\end{quote}

\begin{enumerate}
\def\labelenumi{\arabic{enumi}.}
\setcounter{enumi}{4}
\item
  \textbf{Implementation 34}

  \begin{enumerate}
  \def\labelenumii{\arabic{enumii}.}
  \item
    Environment Setup 34
  \item
    Data Collection and Processing 32
  \item
    Attack Simulation 32
  \item
    Model Development 23
  \item
    System Integration 43
  \end{enumerate}
\item
  \textbf{Output screens and Results 44}
\item
  \textbf{Limitations and Future Enhancements} 68
\end{enumerate}

\begin{quote}
7.1 Limitations

7.2 Future Enhancements 32

7.2.1 Technical Improvements 32

7.2.2 Applications 32
\end{quote}

\begin{enumerate}
\def\labelenumi{\arabic{enumi}.}
\setcounter{enumi}{7}
\item
  \textbf{Conclusion 69}
\item
  \textbf{Appendix A -- Keywords 33}
\item
  \textbf{References 32}
\end{enumerate}

\textbf{LIST OF FIGURES}

4.1 IoT Setup (ESP32 + DHT11 Sensors + MQTT Broker)
\ldots\ldots\ldots\ldots\ldots\ldots\ldots\ldots\ldots\ldots\ldots\ldots.
32

5.2 ARP Spoofing Attack Setup
\ldots\ldots\ldots\ldots\ldots\ldots\ldots\ldots\ldots\ldots\ldots\ldots\ldots\ldots\ldots\ldots\ldots\ldots\ldots\ldots\ldots\ldots\ldots{}
32

5.5 Machine Learning Model Workflow
\ldots\ldots\ldots\ldots\ldots\ldots\ldots\ldots\ldots\ldots\ldots\ldots\ldots\ldots\ldots\ldots\ldots\ldots..\ldots.
32

5.8 Snort IDS Detection and Alert Flow
\ldots\ldots\ldots\ldots\ldots\ldots\ldots\ldots\ldots\ldots\ldots\ldots\ldots\ldots\ldots\ldots\ldots\ldots\ldots\ldots{}
32

6.1 Real-Time Sensor Monitoring Graph
\ldots\ldots\ldots\ldots\ldots\ldots\ldots\ldots\ldots\ldots\ldots\ldots\ldots\ldots\ldots\ldots\ldots\ldots\ldots...
32

\textbf{LIST OF TABLES}

2.1 Comparative Analysis of Existing vs Proposed System

5.1 Dataset Description

5.2 Feature Set Description

6.1 Performance Evaluation Results

A.1 List of Keywords and Definitions

\textbf{Chapter 1}

\textbf{Introduction}

\textbf{1.1 History}

The evolution of building management systems has progressed from manual
control mechanisms to fully automated smart environments powered by IoT
technologies. Earlier systems relied on human intervention to monitor
and regulate environmental conditions such as temperature and humidity.
While these systems were simple, they lacked efficiency and scalability.

With the advancement of IoT, smart building systems began using
distributed sensors and automated controllers to make real-time
decisions. Technologies such as ESP32-based sensor nodes and MQTT
communication enabled continuous data collection and remote monitoring.
These systems improved operational efficiency and reduced human effort.

However, as automation increased, so did the dependency on sensor data.
Most modern systems focus primarily on securing communication channels
to ensure data Authenticity during Transmission. Despite this, they
assume that the received data is always correct. In reality, sensors can
malfunction or be manipulated, resulting in incorrect but believable
data being used for decision-making.

This limitation led to the need for systems that not only secure data
transmission but also validate the correctness of data. Recent
advancements in cybersecurity and machine learning have enabled the
development of intelligent detection systems capable of identifying
anomalies and malicious behavior in real time.

The proposed system builds upon these developments by combining
IoT-based data collection with hybrid intrusion detection techniques,
ensuring that automated decisions are based on reliable and trustworthy
data.

\textbf{1.2 Requirement Analysis}

The proposed system requires a combination of IoT infrastructure,
network monitoring, and intelligent data analysis to detect anomalies
and potential attacks in real time. The requirements focus on ensuring
accurate data collection, reliable detection mechanisms, and secure
system operation. These requirements are categorized into functional and
non-functional requirements.

\textbf{Functional Requirements}

The system must collect real-time environmental data from multiple IoT
sensors (ESP32 with DHT11).

• The system must transmit sensor data to a central server using MQTT
protocol.

• The system must support real-time monitoring of incoming sensor data.

• The system must detect abnormal patterns in data using rule-based
logic (spike, drift, drop, noise).

• The system must implement machine learning to classify normal and
attack conditions.

• The system must analyze data in temporal sequences rather than
individual readings.

• The system must detect replay attacks by comparing current data
patterns with historical sequences.

• The system must monitor network traffic using an intrusion detection
system (Snort).

• The system must generate alerts for suspicious activities at both
network and data levels.

• The system must integrate multiple detection methods into a unified
decision system.

• The system must provide visualization of sensor data and detected
anomalies.

\textbf{Non-Functional Requirements}

• Security\textbf{:} The system must ensure secure communication and
detect malicious data manipulation.

• Reliability\textbf{:} The system must provide consistent and accurate
detection under continuous operation.

• Scalability\textbf{:} The system should support additional sensors and
nodes in the future.

• Performance\textbf{:} The system must process and analyze data in real
time with minimal delay.

• Accuracy\textbf{:} The system should minimize false positives and
false negatives in detection.

• Usability\textbf{:} The system should provide clear and understandable
output through visualization.

• Robustness\textbf{:} The system should handle noisy, incomplete, or
inconsistent data effectively.

\textbf{1.3 Main Objective}

The main objective of this project is to design and develop a hybrid
intrusion detection system for IoT-based smart building environments
that ensures reliable and secure decision-making. The system aims to
detect anomalies and malicious activities in real time by analyzing both
sensor data and network behavior, preventing incorrect decisions caused
by faulty or manipulated data.

\textbf{1.4 Sub Objectives}

• To collect real-time environmental data using IoT devices such as
ESP32 and DHT11 sensors.

• To establish reliable communication between devices using MQTT
protocol.

• To simulate real-world attacks such as ARP spoofing for generating
realistic datasets.

• To detect abnormal patterns in sensor data using rule-based
techniques.

• To develop a machine learning model for classifying normal and attack
conditions.

• To analyze temporal behavior of data using sequence-based processing.

• To implement replay attack detection using pattern comparison
techniques.

• To monitor network traffic using an intrusion detection system
(Snort).

• To integrate network-level and data-level detection into a unified
system.

• To provide real-time visualization of system behavior and detected
anomalies.

\textbf{1.5 System Workflow}

The system follows a structured workflow to collect, analyze, and detect
anomalies in IoT-based smart building environments. The process
integrates data collection, attack simulation, detection mechanisms, and
final decision-making.

• Data Collection: Environmental data such as temperature and humidity
is collected from multiple DHT11 sensors connected to ESP32 devices.

• Data Transmission: The collected data is transmitted in real time to
an MQTT broker for centralized processing.

• Attack Simulation\textbf{:} A real-world attack scenario is created
using ARP spoofing to intercept and manipulate sensor data during
transmission.

• Data Processing: The incoming data is cleaned and prepared for
analysis, including grouping into sequences for temporal evaluation.

• Feature Extraction: Important statistical and temporal features are
extracted from the data to identify patterns and behavior.

• Detection (Rule-Based): Basic anomalies such as spikes, drift, noise,
and drops are detected using predefined rules.

• Detection (Machine Learning): A trained Random Forest model classifies
the data into normal or different attack categories.

• Replay Detection: The system compares current data patterns with past
sequences to detect repeated or replayed data.

• Network Monitoring: Snort IDS monitors network traffic to identify
suspicious activities such as scanning or unauthorized access.

• Decision Making: All detection outputs are combined in a hybrid
decision system to generate a final result.

• Output and Visualization: The system displays results through
real-time graphs and alerts for detected anomalies.

\textbf{Chapter 2}

\textbf{SYSTEM ANALYSIS}

\textbf{2.1 Introduction}

The increasing adoption of IoT-based systems in smart buildings has
enabled automated monitoring and control of environmental conditions
using sensor data. These systems rely heavily on continuous data streams
from multiple devices to make real-time decisions, improving efficiency
and reducing manual effort.

However, most existing systems focus mainly on securing communication
channels while assuming that the received data is always correct. In
practice, sensor data can be faulty, noisy, or even intentionally
manipulated through cyber-attacks. This leads to situations where
systems function normally but operate on incorrect information,
resulting in poor decisions and potential risks.

To address these challenges, there is a growing need for intelligent
systems that can detect anomalies not only at the data level but also at
the network level. By combining IoT data analysis with cybersecurity
techniques, it becomes possible to identify both abnormal sensor
behavior and suspicious network activities.

This chapter analyzes the current approaches used in IoT systems,
highlights their limitations, and explains the need for a hybrid
detection system that improves reliability and security in smart
building environments.

\textbf{2.2 Existing System}

The rapid growth of Internet of Things (IoT) systems has introduced
significant security challenges due to continuous communication between
connected devices. Intrusion Detection Systems (IDS) have been widely
used to monitor network activities and detect malicious behavior in IoT
environments. IDS plays a crucial role in identifying abnormal patterns
and protecting systems from cyber threats {[}1{]}.

Existing IDS approaches are mainly classified into signature-based and
anomaly-based detection systems. Signature-based systems detect known
attacks using predefined patterns, while anomaly-based systems identify
deviations from normal behavior, making them capable of detecting
unknown threats but often producing higher false positives. Suganthi et
al. discussed the challenges of intrusion detection in IoT due to device
heterogeneity and network complexity \textbf{{[}2{]}.}

Recent research has focused on integrating machine learning techniques
into IDS to improve detection accuracy. Rahman et al. highlighted that
machine learning models can effectively detect complex attack patterns
in IoT networks by learning from historical data \textbf{{[}3{]}.}

Furthermore, Kikissagbe et al. demonstrated that machine learning-based
intrusion detection systems enhance detection performance by adapting to
evolving threats and reducing false alarms \textbf{{[}4{]}.}

Hybrid intrusion detection approaches have also been proposed to combine
multiple detection techniques. Alakhras et al. emphasized that hybrid
IDS models improve scalability and detection accuracy by integrating
different strategies such as rule-based and machine learning methods
\textbf{{[}5{]}.}

Despite these advancements, most existing systems focus either on
network-level monitoring or data-level anomaly detection independently.
This separation limits their effectiveness in real-world IoT
environments where both network behavior and sensor data must be
analyzed together. Additionally, challenges such as real-time
processing, scalability, and handling dynamic IoT environments remain
open research problems \textbf{{[}6{]}}.

These limitations highlight the need for a hybrid and integrated
approach that combines multiple detection techniques to improve the
reliability and security of IoT-based systems.

\textbf{2.3 Research Gap}

The analysis of existing systems reveals that most intrusion detection
approaches focus either on network-level monitoring or data-level
anomaly detection independently. Signature-based systems are limited to
detecting known attacks, while anomaly-based and machine learning
approaches require large datasets and often struggle with real-time
implementation.

Additionally, existing solutions lack integration between sensor data
analysis and network traffic monitoring, which is critical in IoT
environments where both data manipulation and network-based attacks can
occur simultaneously.

Furthermore, limited research addresses real-time attack scenarios such
as ARP spoofing in practical IoT deployments. Many systems are tested on
simulated datasets rather than real-world manipulated data.

To overcome these limitations, there is a need for a hybrid system that
combines rule-based detection, machine learning, and network-level
monitoring into a unified framework. The proposed system addresses this
gap by integrating IoT data analysis with network intrusion detection
and real-time attack simulation, improving both accuracy and
reliability.

\textbf{2.4 Comparative Analysis}

{\def\LTcaptype{none} % do not increment counter
\begin{longtable}[]{@{}
  >{\raggedright\arraybackslash}p{(\linewidth - 8\tabcolsep) * \real{0.2598}}
  >{\raggedright\arraybackslash}p{(\linewidth - 8\tabcolsep) * \real{0.1818}}
  >{\raggedright\arraybackslash}p{(\linewidth - 8\tabcolsep) * \real{0.1429}}
  >{\raggedright\arraybackslash}p{(\linewidth - 8\tabcolsep) * \real{0.1428}}
  >{\raggedright\arraybackslash}p{(\linewidth - 8\tabcolsep) * \real{0.2727}}@{}}
\toprule\noalign{}
\begin{minipage}[b]{\linewidth}\centering
\textbf{Citation}
\end{minipage} & \begin{minipage}[b]{\linewidth}\centering
\textbf{Detection Type}
\end{minipage} & \begin{minipage}[b]{\linewidth}\centering
\textbf{ML Usage}
\end{minipage} & \begin{minipage}[b]{\linewidth}\centering
\textbf{Network Monitoring}
\end{minipage} & \begin{minipage}[b]{\linewidth}\centering
\textbf{Remarks}
\end{minipage} \\
\midrule\noalign{}
\endhead
\bottomrule\noalign{}
\endlastfoot
{[}1{]} Sarhan et al. & Anomaly-Based & Applied & Not Applied & Focus on
feature selection for ML-based IDS \\
{[}2{]} Hindy et al. & ML-Based & Applied & Not Applied & MQTT-based
dataset and attack detection \\
{[}3{]} Vitorino et al. & ML Comparison & Applied & Not Applied &
Comparative study of ML techniques \\
{[}4{]} Ali et al. & Hybrid ML & Applied & Applied & Multi-layer
detection for ARP spoofing \\
{[}5{]} Roesch et al. (Snort) & Signature-Based & Not Applied & Applied
& Strong network traffic monitoring \\
{[}6{]} Traditional IDS & Signature-Based & Not Applied & Applied &
Detects known attacks only \\
\end{longtable}
}

\textbf{Chapter 3}

\textbf{System Design}

\textbf{3.1 Proposed System}

The proposed system is a Hybrid Intrusion Detection System (IDS)
designed for IoT-based smart building environments to ensure secure and
reliable decision-making. The system integrates IoT data collection,
network monitoring, and intelligent anomaly detection into a unified
framework.

The system begins with multiple IoT sensor nodes (ESP32 with DHT11
sensors) deployed in different zones to collect environmental data such
as temperature and humidity. These sensors continuously transmit
real-time data to an MQTT broker, enabling centralized monitoring and
processing.

To simulate real-world attack scenarios, a Man-in-the-Middle (MITM)
attack is performed using ARP spoofing, allowing interception and
manipulation of live sensor data. This approach enables the generation
of realistic datasets, improving the effectiveness of the detection
system.

The proposed system incorporates a multi-layer detection mechanism:

• A rule-based detection layer identifies simple anomalies such as
sudden spikes, gradual drift, noise, and abnormal drops in sensor data.

• A machine learning layer uses a Random Forest model trained on
real-time generated datasets to classify normal and attack conditions.
Machine learning-based IDS have been widely used in IoT environments to
improve detection accuracy and adaptability.

• A replay detection mechanism compares current data sequences with
historical patterns to detect repeated or replayed data injections.

• A network monitoring layer using Snort IDS analyzes network traffic to
detect suspicious activities such as scanning, probing, and unauthorized
access. IDS systems play a crucial role in identifying abnormal behavior
in IoT networks.

All detection components are integrated into a hybrid decision system,
which combines outputs from different layers to generate a final
decision with higher accuracy and reduced false positives. Research
shows that hybrid IDS approaches improve detection performance by
combining multiple techniques.

The system also includes a real-time visualization module that provides
graphical representation of sensor data and detected anomalies, enabling
easier monitoring and analysis.

Overall, the proposed system provides a scalable and efficient solution
for detecting both data-level anomalies and network-level attacks,
ensuring that automated decisions are based on reliable and trustworthy
information in IoT environments.

\textbf{3.2 System Architecture}

The proposed system follows a modular and layered architecture that
integrates IoT data collection, network monitoring, and intelligent
anomaly detection. Each component works together to ensure accurate
detection of both data-level and network-level attacks in real time.

The system consists of the following main components:

\textbf{•} IoT Layer:\textbf{\hfill\break
}This layer includes ESP32 devices connected with DHT11 sensors that
continuously collect environmental data such as temperature and humidity
from different zones.

• Communication Layer (MQTT Broker):\\
Sensor data is transmitted in real time to an MQTT broker, which acts as
a central communication hub between IoT devices and the processing
system.

• Attack Simulation Layer:\\
This layer performs ARP spoofing-based Man-in-the-Middle (MITM) attacks
to intercept and manipulate sensor data during transmission, enabling
realistic dataset generation.

• Data Processing Layer:\textbf{\hfill\break
}The collected data is cleaned and organized into sequences for further
analysis. Feature extraction is performed to identify meaningful
patterns in the data.

\textbf{•} Detection Layer:\\
This layer consists of multiple detection mechanisms:

\begin{enumerate}
\def\labelenumi{\arabic{enumi}.}
\item
  Rule-based detection for identifying simple anomalies
\item
  Machine learning model (Random Forest) for classification
\item
  Replay detection for identifying repeated data patterns
\end{enumerate}

• Network Monitoring Layer:\\
A Snort-based IDS monitors network traffic and detects suspicious
activities such as scanning, probing, and unauthorized access.

• Decision Layer:\\
All detection outputs are combined in a hybrid decision system to
produce a final result with improved accuracy and reduced false
positives.

• Visualization Layer:\\
The system provides real-time graphs and outputs to monitor sensor
behavior and detected anomalies.

\textbf{3.3 Modules}

The system is divided into multiple modules, each responsible for a
specific task to ensure efficient detection of anomalies and attacks.

\textbf{•} IoT Data Collection Module:\textbf{\hfill\break
}This module collects real-time environmental data from ESP32 devices
connected with DHT11 sensors.

\textbf{•} Communication Module (MQTT):\\
Handles the transmission of sensor data from IoT devices to the central
system using the MQTT protocol.

\textbf{•} Attack Simulation Module:\\
Implements ARP spoofing to simulate real-world Man-in-the-Middle (MITM)
attacks and generate manipulated data.

\textbf{•} Data Processing Module\textbf{:}\\
Preprocesses incoming data, including cleaning, formatting, and
organizing it into sequences for analysis.

• Feature Extraction Module:\\
Extracts important statistical and temporal features from sensor data to
identify patterns.

\textbf{•} Rule-Based Detection Module:\\
Detects basic anomalies such as spikes, drift, noise, and sudden drops
using predefined rules.

\textbf{•} Machine Learning Module:\\
Uses a Random Forest model to classify data into normal and different
attack categories.

\textbf{•} Replay Detection Module:\\
Identifies repeated or replayed data patterns by comparing current
sequences with historical data.

\textbf{•} Network Monitoring Module (Snort IDS):\\
Monitors network traffic and detects suspicious activities such as
scanning and unauthorized access.

\textbf{•} Decision Module (Hybrid System):\\
Combines outputs from all detection layers to generate a final decision
with improved accuracy.

\textbf{•} Visualization Module:\\
Displays real-time graphs and outputs for monitoring sensor data and
detected anomalies.

\textbf{3.4 Design Diagrams}

3.4.1 Use Case Diagram

Figure 3.1:

3.4.2 Activity Diagram

Figure 3.2:

\textbf{Chapter 4}

\textbf{Technology Stack and System}

\textbf{Specifications}

\subsubsection{\texorpdfstring{\textbf{4.1
Overview}}{4.1 Overview}}\label{overview}

The proposed system is built using a combination of IoT hardware,
communication protocols, and intelligent software components. It
integrates real-time data collection, network monitoring, and anomaly
detection into a unified framework. The system is designed to be
lightweight, scalable, and suitable for IoT environments.

\subsubsection{\texorpdfstring{\textbf{4.2 Hardware
Components}}{4.2 Hardware Components}}\label{hardware-components}

\textbf{•} ESP32 Microcontroller:\\
ESP32 is a low-cost, energy-efficient microcontroller with built-in
Wi-Fi and Bluetooth capabilities, widely used in IoT systems for data
collection and communication

\textbf{•} DHT11 Sensors:\\
Used for collecting environmental data such as temperature and humidity
from different zones.

\textbf{•} Network Environment (Wi-Fi):\\
Provides connectivity between IoT devices and the central system for
real-time communication.

\subsubsection{\texorpdfstring{\textbf{4.3 Software
Components}}{4.3 Software Components}}\label{software-components}

\textbf{•} Programming Language: Python

\textbf{•} Libraries: NumPy, Pandas, Scikit-learn, Matplotlib

\textbf{•} MQTT Broker (Eclipse Mosquitto):\\
Used for lightweight and efficient communication between IoT devices and
the processing system.

\textbf{•} Snort IDS:\\
An open-source network intrusion detection system used to monitor
network traffic and detect suspicious activities such as scanning and
unauthorized access

\textbf{•} Attack Simulation Tools\textbf{:} Bettercap, Wireshark

\subsubsection{\texorpdfstring{\textbf{4.4 System
Configuration}}{4.4 System Configuration}}\label{system-configuration}

• Multi-sensor IoT setup (3 ESP32 + DHT11 nodes)

• Real-time data transmission using MQTT protocol

• Dataset generated through real-time attack simulation (ARP spoofing)

• Machine learning model: Random Forest Classifier

• Feature engineering using statistical and temporal features

• Window-based data processing for sequence analysis.

\textbf{4.5 Design Considerations}

\textbf{4.5.1 Scalability}

The system is designed to support additional IoT devices and sensors
without significant changes to the architecture.

\textbf{4.5.2 Reliability}

Multiple detection layers (rule-based, ML, and network monitoring)
ensure consistent and accurate results even under attack conditions.

\textbf{4.5.3 Security}

The system addresses both network-level and data-level threats,
including ARP spoofing and data manipulation attacks.

\textbf{Chapter 5}

\textbf{Implementation}

\textbf{5.1 Environment Setup}

The system was developed using a combination of hardware and software
tools. ESP32 microcontrollers were connected with DHT11 sensors to
collect environmental data. All devices were connected over a common
Wi-Fi network.

The software environment included Python for data processing and machine
learning, Eclipse Mosquitto as the MQTT broker for communication, and
Snort IDS for network monitoring. Tools such as Bettercap and Wireshark
were used for attack simulation and traffic analysis.

\subsubsection{\texorpdfstring{\textbf{5.2 Data Collection and
Processing}}{5.2 Data Collection and Processing}}\label{data-collection-and-processing}

Real-time environmental data such as temperature and humidity was
collected from multiple sensor nodes. The data was transmitted to the
MQTT broker and then received by a Python-based subscriber system.

The collected data was cleaned and structured into a dataset. It was
further processed into fixed-size windows to capture temporal behavior.
Feature extraction techniques were applied to generate meaningful
attributes such as mean, standard deviation, and rate of change.

\subsubsection{\texorpdfstring{\textbf{5.3 Attack
Simulation}}{5.3 Attack Simulation}}\label{attack-simulation}

A real-world attack scenario was implemented using ARP spoofing. The
attacker system intercepted communication between ESP32 devices and the
MQTT broker.

Using this approach, sensor data was manipulated in real time to
simulate different attack patterns such as spikes, drift, and repeated
data. This helped in generating realistic datasets for training and
testing the detection system.

\subsubsection{\texorpdfstring{\textbf{5.4 Model
Development}}{5.4 Model Development}}\label{model-development}

A Random Forest classifier was developed for anomaly detection. The
model was trained using the generated dataset containing both normal and
attack data.

The system used temporal window-based learning instead of individual
data points. The model classified data into different categories such as
normal behavior, injection, drift, noise, drop, and replay attacks.

Feature scaling and preprocessing were applied to improve model
performance and accuracy.

\subsubsection{\texorpdfstring{\textbf{5.5 System
Integration}}{5.5 System Integration}}\label{system-integration}

All components of the system were integrated into a unified framework.
Sensor data collection, attack simulation, detection mechanisms, and
network monitoring were combined into a hybrid detection pipeline.

The system processed incoming data in real time, applied rule-based and
machine learning detection, and correlated results with network-level
alerts from Snort ID.

\textbf{Chapter 6}

\textbf{Output Screens and Results}

Figure 6.1:

\textbf{Chapter 7}

\textbf{Limitations and Future Enhancements}

\textbf{7.1 Limitations}

The system currently operates in a controlled environment, and
real-world deployment may introduce additional challenges such as
network instability and device heterogeneity.

• The machine learning model performance depends on the quality and
diversity of the dataset. Limited variations in attack scenarios may
affect generalization.

• Real-time processing may introduce slight delays when handling large
volumes of data or multiple sensor nodes simultaneously.

• The system relies on predefined rules for certain detections, which
may not cover all possible attack patterns.

• Snort IDS is used in detection mode only and does not actively block
malicious traffic.

\textbf{7.2 Future Enhancements}

\textbf{7.2.1 Technical Improvements}

• Integration of deep learning models (such as LSTM or CNN) to improve
detection of complex temporal patterns.

• Implementation of automated response mechanisms (IPS) to block
malicious activities in real time.

• Expansion of the dataset with more real-world attack scenarios to
improve model robustness.

• Optimization of the system for large-scale deployments with multiple
IoT devices.

• Development of a dashboard-based interface for advanced monitoring and
control.

\textbf{7.2.2 Applications}

\textbf{•} Smart Buildings\textbf{:} Enhancing reliability and security
of automated HVAC and environmental systems.

\textbf{•} Industrial IoT (IIoT): Monitoring and protecting industrial
sensor networks from cyber threats.

• Healthcare Systems: Ensuring secure monitoring of medical IoT devices
and patient data.

• Smart Cities: Securing large-scale IoT infrastructures used in
traffic, energy, and surveillance systems.

• Research and Cybersecurity Testing: Providing a platform for studying
IoT attack scenarios and developing advanced detection techniques.

\textbf{Chapter 8}

\textbf{Conclusion}

This project presents a hybrid intrusion detection system designed for
IoT-based smart building environments. The system addresses a critical
issue in modern automated systems, where decisions are made based on
sensor data that may be incorrect or manipulated.

The proposed solution integrates multiple components including IoT-based
data collection, real-time communication using MQTT, network monitoring
through Snort IDS, and intelligent anomaly detection using rule-based
techniques and machine learning. A real-world attack scenario using ARP
spoofing was implemented to generate realistic datasets and evaluate
system performance.

By combining data-level analysis with network-level monitoring, the
system provides a more comprehensive approach to intrusion detection
compared to traditional methods. The use of temporal analysis and a
hybrid decision-making mechanism improves detection accuracy while
reducing false positives.

The results demonstrate that the system is capable of identifying
various types of anomalies and attacks in real time, making it suitable
for practical deployment in IoT environments. Overall, the project
highlights the importance of validating both data security and data
correctness in cyber-physical systems.

The proposed system provides a scalable and effective solution for
improving the reliability and security of smart building systems and can
be extended further to support advanced detection and prevention
mechanisms in the future.

\textbf{Appendix A}

\textbf{Keywords}

{\def\LTcaptype{none} % do not increment counter
\begin{longtable}[]{@{}
  >{\raggedright\arraybackslash}p{(\linewidth - 2\tabcolsep) * \real{0.2647}}
  >{\raggedright\arraybackslash}p{(\linewidth - 2\tabcolsep) * \real{0.7353}}@{}}
\toprule\noalign{}
\begin{minipage}[b]{\linewidth}\centering
\textbf{Keyword}
\end{minipage} & \begin{minipage}[b]{\linewidth}\centering
\textbf{Definition}
\end{minipage} \\
\midrule\noalign{}
\endhead
\bottomrule\noalign{}
\endlastfoot
IoT (Internet of Things) & A network of interconnected devices that
collect and exchange data in real time. \\
Intrusion Detection System (IDS) & A system that monitors network or
system activities to detect malicious behavior or security threats. \\
MQTT Protocol & A lightweight messaging protocol used for communication
between IoT devices and servers. \\
ESP32 & A low-cost microcontroller with built-in Wi-Fi used for
IoT-based data collection and communication. \\
DHT11 Sensor & A sensor used to measure temperature and humidity in IoT
environments. \\
ARP Spoofing & A type of cyber-attack where an attacker intercepts
communication by manipulating network address mappings. \\
Machine Learning & A method that enables systems to learn patterns from
data and make predictions or classifications. \\
Random Forest & A machine learning algorithm that uses multiple decision
trees for accurate classification. \\
Anomaly Detection & The process of identifying unusual patterns that do
not match expected system behavior. \\
Rule-Based Detection & A method that uses predefined rules to identify
specific types of anomalies or attacks. \\
Replay Attack & An attack where previously captured data is resent to
manipulate system behavior. \\
Snort IDS & An open-source network intrusion detection system used to
monitor and analyze network traffic. \\
Feature Engineering & The process of extracting meaningful features from
raw data for analysis and model training. \\
Temporal Analysis & The analysis of data over time to detect patterns
and changes in behavior. \\
Hybrid Detection System & A system that combines multiple detection
techniques such as rule-based, machine learning, and network
monitoring. \\
\end{longtable}
}

\textbf{References}

{[}1{]} M. Berhili et al., ``Intrusion Detection Systems in IoT Based on
Machine Learning: A Review,'' \emph{Procedia Computer Science}, 2024.

{[}2{]} A. Kumar and R. Singh, ``Machine Learning for IoT Intrusion
Detection: A Comprehensive Survey,'' \emph{Journal of Network and
Computer Applications}, vol. 222, 2023.

{[}3{]} S. Yaras et al., ``IoT-Based Intrusion Detection System Using
Deep Learning Techniques,'' \emph{Electronics (MDPI)}, vol. 13, 2024.

{[}4{]} A. Almotairi et al., ``Enhancing Intrusion Detection in IoT
Using Machine Learning-Based Models,'' \emph{Journal of Information
Security}, 2024.

{[}5{]} S. Li et al., ``Explainable AI-Based Intrusion Detection in IoT
Systems,'' \emph{ScienceDirect}, 2025.

{[}6{]} S. Fraihat et al., ``Advanced Intrusion Detection for IoT
Devices Using Federated Learning,'' \emph{Springer}, 2025.

{[}7{]} J. Ashraf et al., ``Real-Time Intrusion Detection System for IoT
Networks,'' \emph{Applied Sciences}, 2025.

{[}8{]} M. Ali et al., ``Intrusion Detection in IoT Networks for MitM
Attack Detection Using ML and DL,'' \emph{Springer}, 2025.
